\section{Open problems worked}\label{cx:sec:worked}

Notation as in the introduction of this chapter, with $\gamma_\rho$ the height of the zero $\rho=\sigma_\rho+i\gamma_\rho$, which is the book's ordinate of $\rho$ only under the Riemann hypothesis. The zero free half plane $\Re s>\tfrac78$ is announced, formalised in Lean and not refereed \cite{OAI26a,OAI26b,Liu26a}, as described in Section~\ref{lv:sec:qrh}: every statement below that uses it is conditional on it.

Three open problems were posed for this chapter: the rate of EM at the boundary (answered by experiment, observed and not proved), the value of the latent Brownian times in the sampler (a comparison and an experiment), and a fast algorithm for $u_1$ (answered at fixed accuracy). The global minimum of Section~\ref{cx:sec:limits} is extended to every level in $[\tfrac78,1]$ and then to every level under the matching hypothesis, with the slope of the excess, and the a posteriori bound for the backward flow is worked up to a certificate in ball arithmetic at the levels $\tfrac12$ and $\tfrac34$.

\subsection{The rate at the boundary}\label{cx:sec:rateworked}

The question was whether the factor of about $15$ by which acceleration beat EM at $\theta'=\tfrac12$ is generic, and how the constants depend on the grid. The runs show a grid independent constant, an EM exponent close to $1$ for the first few thousand steps that then rises, and an accelerated scheme whose restarts make the objective a sawtooth with no single exponent; all observed, not proved. The runs use the conventions of Section~\ref{cx:sec:accel}.

Grid independence. The full range problem at $\theta'=\tfrac12$ was run for $3000$ steps at $h=0.04$, $0.02$ and $0.01$. With $R^*$ the smallest residual reached by either method, $k(R_k-R^*)$ for EM is the same on the three grids to four digits (the values of Section~\ref{cx:sec:accel}), and the ratio of EM steps to accelerated steps needed to reach EM's residual at $3000$ steps is $23.1$ ($3000/130$) on all three. We read the iteration as a discretisation of the continuous Richardson and Lucy iteration for measures, whose iterates stay resolved by the grid over thousands of steps. The bound $D(x^*\|x^{(0)})$ of Theorem~\ref{cx:thm:emrate}, which grows like $\log(1/h)$ at an atomic target, is therefore not the active constraint, and the constant of order $50$ appears to belong to the continuous problem.

The exponent. A reference value $R^*$ taken from a finite run biases every exponent estimated from $R_k-R^*$, so we use the three point exponent
\[
p(k)=\log_2\frac{R_k-R_{2k}}{R_{2k}-R_{4k}},
\]
which equals $p$ exactly when $R_k-R^*=Ck^{-p}$ and involves no reference value. For EM it is $1.080$, $1.067$, $1.062$, $1.054$ at $k=500$, $1000$, $2000$, $4000$ at $h=0.08$, equal to about $10^{-3}$ on the grids $0.04$ and $0.02$. After about $4000$ steps it rises, to $1.10$ at $8000$, $1.22$ at $16{,}000$, $1.53$ at $32{,}000$ and $1.84$ at $50{,}000$ in a run of $200{,}000$ steps at $h=0.08$. The residual itself is $9.0\times10^{-4}$, $5.8\times10^{-4}$ and $5.0\times10^{-4}$ at $5\times10^4$, $10^5$ and $2\times10^5$ steps. A geometric rate $R_k-R^*\approx Ce^{-k/K}$ is not claimed: no estimator of $K$ is fixed, and the estimate from three equally spaced differences gives $K\approx32{,}000$ at $k=5\times10^4$ only. The interpretation consistent with the theory of EM on a cone is a sublinear phase at rate $1/k$, while the mass away from the atoms is driven to zero, followed by the local linear convergence of Theorem~\ref{cx:thm:linear} on the identified support. No exponent is claimed to be generic.

Acceleration. For the logarithmic domain scheme with function restart the three point exponent reads $2.8$, $0.6$, $1.2$, $1.2$, $1.3$ at $k=500$, $1000$, $2000$, $4000$, $5000$, a sawtooth whose single values depend on the restart instants; its running value at $k=2000$, $3000$, $5000$ is $1.2$, $1.4$, $1.3$. The speedup over EM, the ratio of step counts to a common residual at $h=0.04$, is $13.5$, $18.9$, $29.6$, $41.7$ and $39.1$ at the residuals EM reaches after $1000$, $2000$, $5000$, $10{,}000$ and $20{,}000$ steps: it grows and saturates near $40$. Five restart rules were compared at $h=0.04$: function restart, no restart, a gradient restart (here the step against the momentum in the logarithmic variables, an adaptation of the rule of O'Donoghue and Cand\`es \cite{OC15}), a fixed restart every $500$ steps, and a momentum capped at $0.9$. They reach EM's residual at $20{,}000$ steps in $512$, $344$, $344$, $344$ and $1991$ steps, with final values $2.0$, $1.7$, $1.5$, $2.8$ and $18\times10^{-5}$ after $20{,}000$ steps. The capped momentum, a heavy ball iteration, is about four times slower than function restart. The exponent without restart is erratic (between $0.2$ and $0.3$ at $k=4000$ and $5000$, not monotone earlier). No restart rule is uniformly best: function restart is the safe one, never increasing $f$, at about one and a half times the steps of the others. There is no generic speedup factor, and the accelerated scheme is the one to run.

\subsection{The latent times do not help the sampler}\label{cx:sec:latent}

The question was whether the three block sampler of Section~\ref{cx:sec:mixture} mixes better than the two block sampler near $\theta'=\tfrac12$. It does not, and for this scan there is a general reason: the two block sampler is the three block sampler with the latent times integrated out analytically. Let $Z$ be the allocations, $X$ the masses and $T$ the times, with $X$ and $T$ independent given $Z$. The allocation marginal of the three block scan is data augmentation through $(X,T)$, that of the collapsed scan is data augmentation through $X$, and for centred $g(Z)$ their reversible kernels satisfy, by conditional Jensen,
\[
\langle g,K_3g\rangle=\E\bigl[\E(g\mid X,T)^2\bigr]\ \ge\ \E\bigl[\E(g\mid X)^2\bigr]=\langle g,K_2g\rangle ,
\]
so the additional augmentation cannot improve the $L^2$ spectral contraction of the allocation marginal. This is the comparison of Liu, Wong and Kong \cite{LiuWongKong94}, in the scope proved here: it orders neither wall clock efficiency nor finite run autocorrelations for other scans and observables.

Both samplers were run on real data (the author's runs, not repeated). The level one density on the window $100\le v\le120$, sampled at spacing $h=0.05$, was turned into Poisson pseudo counts with means
\[
200\,h\,u_1(v_i).
\]
The unknown was on $[98,122]$, the prior $\mathrm{Gamma}(1,0.01)$ per cell, the start the EM solution, and each run had $500$ burn in and $3000$ recorded sweeps. Integrated autocorrelation times were estimated by the initial positive sequence of Geyer \cite{Geyer92}, for the mass within $0.5$ of a zero, the mass in the trough between two zeros, and the log likelihood.

\begin{center}
\small
\begin{tabular}{cccccc}
\toprule
$\theta'$ & sampler & near a zero & trough & log likelihood & seconds\\
\midrule
0.75 & two block & 4.3 & 7.1 & 1.8 & 22\\
0.75 & three block & 4.9 & 9.2 & 2.9 & 65\\
0.55 & two block & 5.4 & 7.8 & 2.5 & 23\\
0.55 & three block & 6.4 & 10.4 & 2.5 & 59\\
\bottomrule
\end{tabular}
\end{center}

\noindent The latent times lengthen the autocorrelation of the two masses by $14$ to $19$ per cent (near a zero) and $30$ to $33$ per cent (trough), and that of the log likelihood by $61$ per cent at $\theta'=0.75$ and not at all at $\theta'=0.55$; they multiply the cost per sweep by $2.6$ to $3$ at both levels, including the one near $\tfrac12$ where they were expected to help. The posterior means agree ($0.80$ for the mass near the zero). The scale mixture keeps its value as an interpretation and a source of proposals, not as a device for speed.

\subsection{The global minimum at every level between $7/8$ and $1$}\label{cx:sec:uniform}

The next problem asked for Theorem~\ref{cx:thm:globalmin}(ii) at every level $\theta\in[\tfrac78,1]$. On the middle range a certificate uniform in $\theta$ seemed to need a two dimensional search; a monotonicity argument removes the second dimension.

\begin{theorem}[global minimum, uniform in the level; conditional]\label{cx:thm:uniform}
Suppose $\zeta$ has no zeros in $\Re s>\tfrac78$. Then for every $\theta\in[\tfrac78,1]$ and every real $v\ne0$ that is not the height of a zero on $\Re s=\theta$,
\[
u_\theta(v)\ >\ u_\theta(0)=\frac1\pi\frac{\xi'}{\xi}(\theta),
\]
and the singular part of $\mathfrak U_\theta$ consists of positive atoms, which occur only if $\theta=\tfrac78$ and a zero lies on $\Re s=\tfrac78$; so $\mathfrak U_\theta-u_\theta(0)\,dv$ is a positive measure. Unconditionally, $\theta\mapsto u_\theta(0)$ is increasing on $[\tfrac12,\infty)$.
\end{theorem}

\begin{proof}
Write $g(\theta,v)=u_\theta(v)-u_\theta(0)$ and $a_\rho=\theta-\sigma_\rho\in[\theta-\tfrac78,\theta-\tfrac18]\subset[0,\tfrac78]$, which uses the hypothesis and the functional equation. Monotonicity of $u_\theta(0)$: $\xi(\tfrac12+x)$ is a positive multiple of $\int_0^\infty\Phi(y)\cosh(xy)\,dy$ with the positive density $\Phi$ of \eqref{eq:ch2:Phi}, and the logarithm of the Laplace transform of a positive measure is strictly convex (H\"older), so $(\xi'/\xi)(\tfrac12+x)$ is strictly increasing in $x$. In particular $u_\theta(0)\le u_1(0)=0.00736$ on the range.

$|v|\le13.5$: every term $P_{a_\rho}(v-\gamma_\rho)$ with $a_\rho>0$ is convex, since $|v-\gamma_\rho|\ge0.63>\tfrac78/\sqrt3$; the sum is strictly convex and even, with its minimum at $0$. This is the argument of Theorem~\ref{cx:thm:globalmin}, valid for all $\theta\le1$ at once.

$|v|\ge1418.9$: Lemma~\ref{cx:lem:pairing} and the zero count of \cite{HSW22} give at least seven ordinates in $(v-8,v+8]$ (the explicit bound is $8.76786$ at $v=1418.9$, so there are at least nine, and it increases with $v$, the derivative of its main term $(2\pi)^{-1}\log\frac{v+8}{v-8}\ge8/(\pi v)$ dominating the derivatives of the error terms, which are $O(1/v)$; this was checked numerically in the computation of Section~\ref{cx:sec:uniform}). Seven suffice: $u_\theta(v)\ge\tfrac72P_{2\theta-1}(8)\ge\tfrac72P_{3/4}(8)=0.01294>u_1(0)\ge u_\theta(0)$, using that $P_c(8)$ increases in $c$ on $[\tfrac34,1]$.

$13.5\le|v|\le1418.9$: by Theorem~\ref{cx:thm:certified} every zero with $|\gamma|<1419.9$ lies on the critical line and is one of the $1000$ certified ordinates (and their negatives), so its width is exactly $a=\theta-\tfrac12\le\tfrac12$, and the zeros with $|\gamma|\ge1419.9$ are at distance at least $1$ from $v$. Two cases.

If some certified zero has $|v-\gamma|\le1.05$, then $u_\theta(v)\ge P_{\theta-1/2}(1.05)\ge P_{3/8}(1.05)=0.0960>u_\theta(0)$, since $P_c(1.05)$ increases in $c$ on $[\tfrac38,\tfrac12]$.

If every certified zero has $|v-\gamma|\ge0.95$, differentiate in $\theta$:
\[
\partial_\theta u_\theta(v)=\frac1\pi\sum_\rho f\bigl(|v-\gamma_\rho|,a_\rho\bigr),\qquad f(d,a)=\frac{d^2-a^2}{(d^2+a^2)^2}.
\]
For the far zeros, $d\ge1>\tfrac78\ge a$, so their terms are positive. For the certified zeros, $a\le\tfrac12$, and $f$ is decreasing in $a$ for $d>a/\sqrt3$ and decreasing in $d$ for $d>\sqrt3\,a$, both satisfied when $d\ge0.95$ and $a\le\tfrac12$. Hence
\[
\partial_\theta u_\theta(v)\ \ge\ \frac1\pi\sum_{|\gamma|<1419.9}f\bigl(|v-\gamma|,\tfrac12\bigr)=:S(v),
\]
a function of the certified ordinates alone, independent of $\theta$. On the other side, $\partial_\theta u_\theta(0)=\pi^{-1}\sum_\rho(\gamma_\rho^2-a_\rho^2)/(a_\rho^2+\gamma_\rho^2)^2\le\pi^{-1}\sum_\rho\gamma_\rho^{-2}$. The sum over the certified ordinates, both signs, is at most $0.044804$, and the sum over $|\gamma|\ge1419.9$ is at most $0.001442$ (the computed value is $0.0014417$) by partial summation against the bound of \cite{HSW22} for $N(T)$, so $\partial_\theta u_\theta(0)\le c^*=0.0147225$ for every $\theta$. The ball arithmetic computation described below certifies $S(v)>c^*$ on every cell of the second case. Therefore $\partial_\theta g\ge0$ on $[\tfrac78,1]$, and $g(\theta,v)\ge g(\tfrac78,v)>0$ by Theorem~\ref{cx:thm:globalmin}(ii).

Atoms: a zero contributes an atom to $\mathfrak U_\theta$ only if its width $a_\rho=0$, that is $\sigma_\rho=\theta$, and $\sigma_\rho\le\tfrac78\le\theta$ forces $\theta=\tfrac78$.
\end{proof}

The computation partitions $[13.5,1418.9]$ into cells of radius $0.1$, bisected where neither case applies to the whole cell. The ordinates enter as the certified sign change enclosures. A cell is of the first kind if for one certified ordinate the largest distance between a point of the cell and a point of its enclosure is at most $1.05$, and of the second kind if the smallest distance between the cell and every enclosure is at least $0.95$; then $S$ is bounded below by evaluating each $f(\cdot,\tfrac12)$ at the largest such distance, since $f$ decreases in $d$. Of the $7{,}245$ cells, $6{,}733$ are of the first kind and $512$ of the second, every cell is certified, and the least certified value of $S$ on a cell of the second kind is $0.066$ against $c^*$.

\begin{corollary}\label{cx:cor:uniform}
Given the zero free half plane $\Re s>\tfrac78$, for every $\theta\in[\tfrac78,1]$ the Thorin measure $\mathfrak U_\theta$ dominates the multiple $u_\theta(0)\,dv$ of Lebesgue measure, with $u_\theta(0)$ increasing from $0.0055147$ at $\theta=\tfrac78$ to $0.0073516$ at $\theta=1$.
\end{corollary}

\subsection{Below $7/8$: a level by level theorem}\label{cx:sec:lbl}

The case $\theta=1$ is Proposition~\ref{bp:prop:v1min} and the case $\theta\downarrow\tfrac12$ is Theorem~\ref{thm:dbeta:minall}(vi).

\begin{theorem}[level by level]\label{cx:thm:lbl}
Let $\tfrac12<\theta\le1$. The following are equivalent.
\begin{enumerate}
\item[(i)] $\zeta$ has no zeros in $\Re s>\theta$.
\item[(ii)] For every $\theta'\in[\theta,1]$ the function $u_{\theta'}$ attains its global minimum at the origin: $u_{\theta'}(v)\ge u_{\theta'}(0)$ for all real $v$ that are not heights of zeros on $\Re s=\theta'$.
\end{enumerate}
When they hold, $\mathfrak U_{\theta'}\ge u_{\theta'}(0)\,dv$ as measures for $\theta'\in[\theta,1]$, and for every $\theta'\in[\theta,1]$, with $c=2\theta'-1$ and the sum over all zeros,
\begin{equation}\label{cx:eq:excess}
u_{\theta'}(v)-u_{\theta'}(0)\ \ge\ \frac{c}{2\pi}\Bigl(\sum_\rho\frac1{c^2+(v-\gamma_\rho)^2}-\pi w^*\Bigr),
\end{equation}
with $w^*=0.0294450$, and the bracket is at least $\pi(0.06174-w^*)>0.1014$ for $13.5\le|v|\le1418.9$ and at least $7/65-\pi w^*>0.0151$ for $|v|\ge1418.9$.
\end{theorem}

\begin{proof}
(ii) implies (i) exactly as in Theorem~\ref{cx:thm:rhmin}: a zero $\rho_0$ of multiplicity $m$ with $\sigma_0>\theta$ makes $u_{\theta'}(\gamma_0)=-m/\varepsilon+O(1)$ at the level $\theta'=\sigma_0-\varepsilon\in[\theta,1]$, which is negative, while $u_{\theta'}(0)=\pi^{-1}(\xi'/\xi)(\theta')>0$.

(i) implies (ii). Fix $\theta'\in[\theta,1]$, write $c=2\theta'-1\in(0,1]$, and let $a_\rho=\theta'-\sigma_\rho\ge0$. Two bounds, one on each side.

At the origin, grouping the zeros into orbits under $\rho\mapsto1-\bar\rho$, the widths in an orbit sum to $c$ (a pair) or equal $c/2$ (a zero on the line), so
\[
u_{\theta'}(0)=\frac1\pi\sum_\rho\frac{a_\rho}{a_\rho^2+\gamma_\rho^2}\ \le\ \frac1\pi\sum_{\text{orbits}}\frac{c}{\gamma^2}\ \le\ \frac c\pi\sum_\rho\frac1{\gamma_\rho^2}\ \le\ c\,c^*,
\]
with $c^*=0.0147225$, the certified value of Theorem~\ref{cx:thm:uniform}: the certified ordinates below $1419.9$ contribute at most $0.044804$ to $\sum_\rho\gamma_\rho^{-2}$ and the rest at most $0.001442$ by \cite{HSW22}. (An orbit contributes at most $\tfrac c2\sum\gamma^{-2}$ over its zeros, which would halve the bound; the weaker form is what the printed constants use.)

Away from the origin, Lemma~\ref{cx:lem:pairing} gives $u_{\theta'}(v)\ge\tfrac12\sum_\rho P_c(v-\gamma_\rho)=\frac{c}{2\pi}\sum_\rho\frac1{c^2+(v-\gamma_\rho)^2}\ge\frac{c}{2\pi}\sum_\rho\frac1{1+(v-\gamma_\rho)^2}$, since $c\le1$. So \eqref{cx:eq:excess} holds with $w^*=2c^*$, and (ii) at $v$ follows from
\begin{equation}\label{cx:eq:W}
W(v)=\frac1\pi\sum_\rho\frac1{1+(v-\gamma_\rho)^2}\ \ge\ 2c^*=w^*,
\end{equation}
a condition on the ordinates alone, with no $\theta'$ in it. Dropping the zeros above height $1419.9$ and those of negative height only lowers $W$. Ball arithmetic certifies \eqref{cx:eq:W} on $[13.5,1418.9]$ with the $1000$ certified positive ordinates, $7027$ cells of radius $0.1$, each term evaluated at the farthest point of the cell; the smallest certified cell bound is $W\ge0.06174$, in the cell about $v=17.4$, between the first two zeros, and the bracket is at least $\pi(0.06174-w^*)=0.10146$. With the negative ordinates included the same cells give $0.0653$. For $|v|\ge1418.9$ the seven ordinates within distance $8$ give $W\ge7/(65\pi)>0.03427>w^*$ and the bracket is at least $7/65-\pi w^*=0.01518$ (with the nine ordinates of the count of \cite{HSW22} it is $9/65-\pi w^*=0.0459$). For $|v|\le13.5$ every term $P_{a_\rho}(v-\gamma_\rho)$ is convex, because $a_\rho\le\theta'-(1-\theta)\le1$ and $|v-\gamma_\rho|\ge0.63>1/\sqrt3$, so $u_{\theta'}$ is strictly convex and even there and its minimum is at $0$. The excess \eqref{cx:eq:excess} is the difference of the two bounds. The positivity of the atoms gives the measure statement.
\end{proof}

Theorem~\ref{cx:thm:lbl} contains the weak form of Theorem~\ref{cx:thm:globalmin}(i) ($\theta=1$), Theorem~\ref{cx:thm:uniform} ($\theta=\tfrac78$) and Theorem~\ref{cx:thm:rhmin} ($\theta\downarrow\tfrac12$). Its proof is shorter because the pairing lower bound and the orbit upper bound both scale linearly in $c$, so the level drops out and one certified function of the ordinates, $W$, decides every level; the monotonicity in $\theta$ is not used.

As $\theta'\downarrow\tfrac12$ the excess over the value at the origin has a lower bound linear in $c=2\theta'-1$, with slope at least $0.1014/(2\pi)=0.0161$ on the certified range and $0.0151/(2\pi)=0.0024$ beyond it. The slope is not sharp, since the pairing lemma halves the density and $W$ was bounded with $c=1$, but the order is right: under the Riemann hypothesis, at the midpoint between two consecutive zeros with gap $G$, $u_{\theta'}(v)=\sum P_{c/2}(v-\gamma)\approx\frac{c}{2\pi}\cdot\frac{8}{G^2}$ for small $c$, and $u_{\theta'}(0)\approx c\,c^*/2$, so the excess is of order $c$ at fixed $v$ (of order $1/c$ at a zero).

\subsection{The slope of the excess}\label{cx:sec:slope}

The losses in \eqref{cx:eq:excess} are avoidable for the zeros whose position is known. By Theorem~\ref{cx:thm:certified} every zero with $|\gamma|<1419.9$ lies on the critical line, so at level $\theta'$ its width is exactly $c/2$ and its contribution to $u_{\theta'}$ is exactly $P_{c/2}(v-\gamma)$, with no pairing. Only the zeros above height $1419.9$, whose real parts are unknown, need Lemma~\ref{cx:lem:pairing} and the orbit bound. Writing $V_c(v)=\sum_{|\gamma|<1419.9}\bigl(\tfrac{c^2}4+(v-\gamma)^2\bigr)^{-1}$, both signs of $\gamma$ included, the two sides become
\begin{gather*}
u_{\theta'}(v)\ \ge\ \frac c{2\pi}\Bigl(V_c(v)+\sum_{|\gamma|\ge1419.9}\frac1{c^2+(v-\gamma)^2}\Bigr),\\
u_{\theta'}(0)\ \le\ \frac c{2\pi}\Bigl(V_c(0)+2\sum_{|\gamma|\ge1419.9}\frac1{\gamma^2}\Bigr),
\end{gather*}
and dropping the positive far sum on the left,
\begin{equation}\label{cx:eq:excess2}
u_{\theta'}(v)-u_{\theta'}(0)\ \ge\ \frac c{2\pi}\,B_c(v),\qquad B_c(v)=V_c(v)-V_c(0)-2\!\!\sum_{|\gamma|\ge1419.9}\!\!\frac1{\gamma^2},
\end{equation}
where the last sum is at most $0.001442$ by \cite{HSW22}, so that $B_c(v)\ge V_c(v)-V_c(0)-0.0029$ (the far sum in the display below is written with $c^2$ in place of $\tfrac{c^2}4$, which is the weaker bound). On $13.5\le|v|\le1418.9$ the function $B_c$ decreases in $c$ (its derivative in $c$ is $\tfrac c2\sum_\gamma\bigl[(\tfrac{c^2}4+\gamma^2)^{-2}-(\tfrac{c^2}4+(v-\gamma)^2)^{-2}\bigr]$, in which the first sum is below $\sum\gamma^{-4}<7.5\times10^{-5}$ while the nearest certified ordinate, within $3.45$ of $v$, alone gives more than $(\tfrac14+3.45^2)^{-2}>6.7\times10^{-3}$ to the second; single terms need not decrease), so its value at $c=0$, which involves only $\sum(v-\gamma)^{-2}$ over the certified ordinates, is a lower bound for the slope as $\theta'\downarrow\tfrac12$.

\begin{proposition}[the slope]\label{cx:prop:slope}
Under the hypothesis of Theorem~\ref{cx:thm:lbl}, for every $\theta'\in[\theta,1]$ and $c=2\theta'-1$,
\[
\inf_{13.5\le|v|\le1418.9}B_c(v)\ \ge\ \begin{cases}0.17575&c=1,\\0.17838&c=\tfrac12,\\0.17904&c=\tfrac14,\\0.17927&c\downarrow0,\end{cases}
\]
each infimum being approached at $v=17.49$ to $17.50$, between the first two zeros; and, for $|v|\ge1418.9$,
\[
B_c(v)+\sum_{|\gamma|\ge1419.9}\frac1{c^2+(v-\gamma)^2}\ \ge\ \frac7{c^2+64}-V_c(0)-0.0029\ \ge\ 0.0599,
\]
while $B_c(v)$ itself tends to $-V_c(0)-2\sum_{|\gamma|\ge1419.9}\gamma^{-2}<0$ as $|v|\to\infty$. Under the Riemann hypothesis the limiting bracket
\[
X(v)=\lim_{c\downarrow0}\frac{2\pi}c\bigl(u_{\theta'}(v)-u_{\theta'}(0)\bigr)=\sum_\rho\Bigl(\frac1{(v-\gamma_\rho)^2}-\frac1{\gamma_\rho^2}\Bigr)
\]
has, in the first gap between the zeros $14.1347$ and $21.0220$, a unique minimiser $v^*\in[17.4978,17.4980]$ (below the midpoint $17.5784$ of the gap), and
\[
0.1838985448\ \le\ X(v^*)\ \le\ 0.1838987563 ;
\]
on the rest of $13.5\le|v|\le1418.9$, $V_0(v)-V_0(0)\ge0.184074$ (the author's computation, not rerun).
\end{proposition}

\begin{proof}
First clause, ball arithmetic as for \eqref{cx:eq:W}: the interval $[13.5,1418.9]$ is cut into cells of radius $0.1$, $V_c$ is bounded below on a cell by evaluating each term at the farthest point of the cell from the enclosure of $\gamma$, $V_c(0)$ is bounded above using the lower endpoints of the enclosures, and the cells whose lower bound falls below the best upper bound plus $0.005$ are bisected. The printed values are valid lower bounds; a rerun with bisection to radius $2^{-14}$ and the enclosures of width $0.02$ gives $0.175808$, $0.178430$, $0.179098$, $0.179322$ at $v=17.493$ to $17.498$. Far range: seven ordinates lie within distance $8$ of $v$; each contributes at least $(c^2+64)^{-1}$, through $V_c(v)$ if it is certified and through the far sum of the first display otherwise. The minimum over $c\in[0,1]$ of $7/(c^2+64)-V_c(0)-0.0029$ is $0.0600$ (at $c=1$, using $V_1(0)\le0.044785$), so $0.0599$ holds.

The last statement. Under the Riemann hypothesis $X(v)=[V(v)-V(0)]+F(v)$ with $V=V_0$ over the certified ordinates and $F(v)=\sum_{\gamma\ge1419.9}f_v(\gamma)$, where
\[
f_v(\gamma)=\frac1{(\gamma-v)^2}+\frac1{(\gamma+v)^2}-\frac2{\gamma^2}=\frac{2v^2(3\gamma^2-v^2)}{\gamma^2(\gamma^2-v^2)^2}\ \ge\ 0\qquad(\gamma>v),
\]
so the uncertified zeros only raise $X$, and $X\ge V-V(0)$ on the whole range. Under the Riemann hypothesis $X''(v)=6\sum_\rho(v-\gamma_\rho)^{-4}>0$ between consecutive ordinates, so $X$ is strictly convex on $(14.1347,21.0220)$ and has a unique minimiser $v^*$ there; on $[13.5,14.1347)$ it is increasing, being even and convex on $(-\gamma_1,\gamma_1)$, and $X(13.5)>2$, the nearest ordinate alone contributing $2.48$ while $\sum_\rho\gamma_\rho^{-2}=0.0462$. For $0<v<\gamma$, $\partial_vf_v(\gamma)=4v(3\gamma^2+v^2)/(\gamma^2-v^2)^3>0$, so $F'\ge0$. With the $1000$ certified ordinates (\texttt{python-flint}, $160$ and $320$ bits) the derivative $Q'$ of $Q=V-V(0)$ is enclosed by $Q'(17.4978)\in[-1.0391209,-1.0367081]\times10^{-5}$ and $Q'(17.4980)\in[7.2326504,7.2567776]\times10^{-6}$, including a rigorous bound $0\le F'\le2.41\times10^{-8}$ for the tail (partial summation against \cite{HSW22}). So $X'$ changes sign exactly once in $[17.4978,17.4980]$. For the value, tangent lines of $Q$ at the two endpoints give the lower bound, and $Q(17.4979)$ plus the tail bound $2.11\times10^{-7}$ the upper. The statement for the rest of the range is the cell computation of the first clause for $V-V(0)$ with the enclosures of the author, not rerun: the rerun of the first clause gives $V-V(0)\ge B_0\ge0.17927$ only.
\end{proof}

The certified values for the four levels differ by two per cent, so the slope is essentially independent of the level: the bracket of \eqref{cx:eq:excess} is at least $0.1757$ rather than $0.1014$ on the middle range and $0.0599$ rather than $0.0151$ beyond, a slope per unit $c$ of at least $0.1757/(2\pi)=0.0279$. By the author's computation (the last clause of Proposition~\ref{cx:prop:slope}, not rerun) the infimum over $13.5\le v\le1418.9$ is attained between the first two zeros, at $v^*$, pushed below their midpoint $17.578$ by the zeros above; its value as $\theta'\downarrow\tfrac12$ is $0.18390$ to five digits under the Riemann hypothesis, the uncertified zeros contributing less than $3\times10^{-7}$ because their paired terms $f_v(\gamma)$ are of order $v^2/\gamma^4$. The two neighbours alone give $8/G^2=0.1687$ with $G=6.887$ at the midpoint of the gap ($0.1689$ at $v^*$), the other $1998$ certified ordinates of both signs add $0.060$, and $V(0)$ takes away $0.0448$. The unconditional certificate \eqref{cx:eq:excess2} is weaker than this by the orbit bound at the origin, $2\sum_{\mathrm{far}}\gamma^{-2}\le0.0029$.

\subsection{A posteriori bounds for the backward flow}\label{cx:sec:apost}

The next problem asked for a bound on the distance between the EM output and the Thorin measure in terms of the residual, so that quantities read from the output become certified. The bound is a duality.

\begin{proposition}[a posteriori window bound]\label{cx:prop:apost}
Let $\theta'\in[\tfrac12,1)$, $a=1-\theta'$, and suppose $\zeta$ has no zeros in $\Re s>\theta'$, so that $\mathfrak U_{\theta'}$ exists and $P_a*\mathfrak U_{\theta'}=u_1\,dv$. Let $\mu$ be a finite positive measure, for instance $\mu=\sum_jx_jh\,\delta_{w_j}$ built from an EM iterate, and let $r=P_a*\mu-u_1$ on $\R$. For a test function $\varphi$ with $\int e^{a|\omega|}|\widehat\varphi(\omega)|\,d\omega<\infty$ put $\psi=\mathcal F^{-1}[e^{a|\omega|}\widehat\varphi\,]$, and assume that $\psi\cdot(P_a*\mu+u_1)$ is integrable (this holds for the Gaussian windows below, by the tails of Lemma~\ref{cx:lem:psiclosed}). Then
\[
\Bigl|\int\varphi\,d\mu-\int\varphi\,d\mathfrak U_{\theta'}\Bigr|\ \le\ \int_\R|r(v)|\,|\psi(v)|\,dv .
\]
\end{proposition}

\begin{proof}
$\widehat{P_a}=e^{-a|\omega|}$, so $P_a*\psi=\varphi$. By Fubini, justified by the integrability hypothesis applied to the absolute integrand, $\int\varphi\,d\mu=\int\psi\,(P_a*\mu)\,dv$ and $\int\varphi\,d\mathfrak U_{\theta'}=\int\psi\,u_1\,dv$. Subtract.
\end{proof}

The residual is needed on all of $\R$: inside the data window it is the EM residual, between the data window and the edge of the grid carrying $\mu$ it is computable from $u_1$, and beyond that grid $r=P_a*\mu-u_1$ with $P_a*\mu(v)\le a\|\mu\|/(\pi(a^2+(|v|-W)^2))$ for $\operatorname{supp}\mu\subset[-W,W]$ and $|v|>W$, against the decay of $\psi$, which is $|\psi(v)|\sim a\sigma\sqrt{2\pi}/(\pi(v-c)^2)$ for the Gaussian window $\varphi=e^{-(v-c)^2/2\sigma^2}$, from the kink of $e^{a|\omega|}$ at the origin. The data must be known to certified accuracy on the computation grid: in the author's experience, interpolating $u_1$ from a coarser grid introduces errors of $10^{-3}$, larger than the EM residual, and the bound fails. And the tail must be measured from the window centre, not from the origin. Centroid certificates come from the odd test function $\varphi_1(v)=(v-c)\varphi(v)$: with certified mass $M\pm b_0$ and certified first moment $m_1\pm b_1$, the $\varphi$ weighted centroid of $\mathfrak U_{\theta'}$ lies within $(|m_1|+b_1)/(M-b_0)$ of $c$. An interval for a single zero needs in addition a bound on the other zeros (Corollary~\ref{cx:cor:zeroloc}).

In floating point, with exact data at spacing $0.02$ on $[-400,400]$, the unknown on $[-440,440]$ and $4000$ accelerated steps, the author found the bound to hold for all $198$ windows centred at the zeros in $(5,395)$, for masses and moments, at $\theta'=\tfrac12$ and $\tfrac34$ and $\sigma=0.5$ and $0.3$ (not repeated here); the certified figures below supersede these.

\subsection{The certificate}\label{cx:sec:apostcert}

Two things stood in the way of a certificate: $\psi$ came from a floating point transform, and $\int|r||\psi|$ was a Riemann sum. Both are removed by replacing the Gaussian $\varphi$ with the function the grid actually sees.

\begin{lemma}[the dual function in closed form]\label{cx:lem:psiclosed}
For $\varphi(v)=e^{-(v-c)^2/2\sigma^2}$ and $a>0$, with $q=a+i(v-c)$ and
\[
E(q)=e^{q^2/2\sigma^2}\operatorname{erfc}\Bigl(-\frac q{\sigma\sqrt2}\Bigr),
\]
the dual functions of $\varphi$ and of $\varphi_1(v)=(v-c)\varphi(v)$ are
\[
\psi(v)=\Re E(q),\qquad\psi_1(v)=\Im\bigl(q\,E(q)\bigr)=-\sigma^2\psi'(v).
\]
Both extend to entire functions of $v$: with $q_\pm=a\pm i(z-c)$, $\Psi(z)=\tfrac12[E(q_+)+E(q_-)]$ and $\Psi_1(z)=\frac1{2i}[q_+E(q_+)-q_-E(q_-)]$. For every $x=v-c\ne0$,
\[
|\psi(v)|\le\frac{\sigma\sqrt{2\pi}}\pi\Bigl(\frac a{x^2}+\frac{B_4}{x^4}\Bigr),\qquad
|\psi_1(v)|\le\frac{\sigma^3\sqrt{2\pi}}\pi\Bigl(\frac{2a}{|x|^3}+\frac{B_5}{|x|^5}\Bigr),
\]
where, with $g(\omega)=e^{a\omega-\sigma^2\omega^2/2}$ and $G(\omega)=\omega g(\omega)$, $B_4=|a^3-3\sigma^2a|+\int_0^\infty|g^{(4)}|$ and $B_5=4|a^3-3\sigma^2a|+\int_0^\infty|G^{(5)}|$.
\end{lemma}

\begin{proof}
By definition,
\begin{gather*}
\psi(v)=\frac1{2\pi}\int e^{a|\omega|}\widehat\varphi(\omega)e^{i\omega v}d\omega=\frac{\sigma\sqrt{2\pi}}\pi\int_0^\infty g(\omega)\cos(\omega x)\,d\omega,\\
\int_0^\infty e^{-p\omega^2+q\omega}d\omega=\tfrac12\sqrt{\pi/p}\,e^{q^2/4p}\operatorname{erfc}\bigl(-q/2\sqrt p\bigr),
\end{gather*}
and the second, classical, integral with $p=\sigma^2/2$ and $q=a+ix$ gives the first formula; at $a=0$ it reduces to $\varphi$, since $\operatorname{erfc}$ of an imaginary argument has real part $1$. Since $\overline{E(q)}=E(\bar q)$, the real part is $\tfrac12[E(q_+)+E(q_-)]$ on the real axis, which is entire in $z$. Differentiating, $E'(q)=(q/\sigma^2)E(q)+\sqrt{2/\pi}/\sigma$, so $\psi'(v)=\Re(iE'(q))=-\sigma^{-2}\Im(qE(q))$; and $\widehat{\varphi_1}=-i\sigma^2\omega\widehat\varphi$ gives $\psi_1=-\sigma^2\psi'$. The tails come from integrating by parts:
\begin{gather*}
\int_0^\infty g\cos(\omega x)\,d\omega=-\frac{g'(0)}{x^2}+\frac{g'''(0)}{x^4}+\frac1{x^4}\int_0^\infty g^{(4)}\cos(\omega x)\,d\omega,\\
\int_0^\infty G\sin(\omega x)\,d\omega=-\frac{G''(0)}{x^3}+\frac{G''''(0)}{x^5}+\frac1{x^5}\int_0^\infty G^{(5)}\cos(\omega x)\,d\omega,
\end{gather*}
with $g'(0)=a$, $g'''(0)=a^3-3\sigma^2a$, $G''(0)=2a$ and $G''''(0)=4(a^3-3\sigma^2a)$; the second integral is the one for $\psi_1=\frac{\sigma^3\sqrt{2\pi}}\pi\int_0^\infty G\sin(\omega x)\,d\omega$.
\end{proof}

The integrals $\int|g^{(4)}|$ and $\int|G^{(5)}|$ are enclosed by writing $g^{(k)}=H_kg$ with $H_{k+1}=H_k'+(a-\sigma^2\omega)H_k$, taking absolute values of the coefficients and integrating in balls. The author reports $B_4=38.0$ and $B_5=550.6$ at $a=\tfrac12$, $\sigma=0.3$ (not rerun); generous, but they multiply $x^{-4}$ and $x^{-5}$ and are used only at distances above $440$.

Now fix the grid $v_i=ih$, $|v_i|\le880$, and let $\psi_h=\sum_i\psi(v_i)h\,\delta_{v_i}$ and $\varphi_h=P_a*\psi_h$. For any finite positive $\mu$ with $P_a*\mu$ finite at the grid points, and $U=\mathfrak U_{\theta'}$,
\begin{equation}\label{cx:eq:exactid}
\int\varphi_h\,d\mu-\int\varphi_h\,d\mathfrak U_{\theta'}=\sum_i\psi(v_i)\,h\,r_i,\qquad r_i=(P_a*\mu)(v_i)-u_1(v_i),
\end{equation}
exactly, by Fubini over a finite sum and $P_a*\mathfrak U_{\theta'}=u_1$. The right side is a finite sum of balls. For $\mu=\sum_jx_jh\delta_{w_j}$, $(P_a*\mu)(v_i)$ is a polynomial product, which Arb \cite{Arb} multiplies with rigorous error bounds in a fraction of a second, the $x_j$ being the double precision EM output taken as exact rationals. The identity does not depend on EM having converged. It remains to compare $\varphi_h$ with $\varphi$:
\[
\int\varphi\,d\mu-\int\varphi\,d\mathfrak U_{\theta'}=\sum_i\psi(v_i)hr_i+\int(\varphi-\varphi_h)\,d\mu-\int(\varphi-\varphi_h)\,d\mathfrak U_{\theta'} .
\]
Write $\varphi_h^\infty$ for the same sum over the whole lattice $h\mathbb Z$. Then $\varphi-\varphi_h=(\varphi-\varphi_h^\infty)+\sum_{|v_i|>880}\psi(v_i)hP_a(\cdot-v_i)$, the first term the trapezoid error of an analytic integrand, the second a truncation controlled by the tails of Lemma~\ref{cx:lem:psiclosed}.

\begin{lemma}[trapezoid error]\label{cx:lem:trap}
Let $\Psi$ be the entire extension of $\psi$, $f_w(v)=\Psi(v)P_a(w-v)$, and assume $h\le\pi\sigma^2/(2a)$. For every real $w$,
\[
|\varphi(w)-\varphi_h^\infty(w)|\ \le\ \min\Bigl(\frac{4S_\Psi}{e^{\pi a/h}-1},\ \frac{E_{\mathrm{dec}}}{(w-c)^2}\Bigr),
\]
where $S_\Psi=2e^{9a^2/8\sigma^2}+1$ bounds $|\Psi|$ on the strip $|\Im v|\le a/2$, and
\begin{gather*}
E_{\mathrm{dec}}=\frac{\sigma\sqrt{2\pi}}{2\pi}\,(4a+J)\,\frac{2e^{-2\pi a/h}}{1-e^{-2\pi a/h}},\\
J=\int_{\R}\bigl((2a+\sigma^2|\omega|)^2+\sigma^2\bigr)e^{2a|\omega|-\sigma^2\omega^2/2}\,d\omega ;
\end{gather*}
at $w=c$ the second term of the minimum is $+\infty$. The same holds for $\psi_1$ and $\varphi_1$ with $b=\tfrac{3a}2$, $S_1=\sigma^3\sqrt{2\pi}\,\pi^{-1}\int_0^\infty te^{bt-\sigma^2t^2/2}dt$ in place of $S_\Psi$ with numerator $4S_1$, $C_1=\sup_{t\ge0}te^{2at-\sigma^2t^2/2}$, $J_1=\int_{\R}\bigl[2(2a+\sigma^2|\omega|)+|\omega|\bigl((2a+\sigma^2|\omega|)^2+\sigma^2\bigr)\bigr]e^{2a|\omega|-\sigma^2\omega^2/2}d\omega$, and
\[
E_{\mathrm{dec},1}=\frac{\sigma^3\sqrt{2\pi}}{2\pi}\,(2aC_1+J_1)\,\frac{2e^{-2\pi a/h}}{1-e^{-2\pi a/h}} .
\]
\end{lemma}

\begin{proof}
The uniform bound is the theorem of Trefethen and Weideman \cite{TW14} for the trapezoid rule on $\R$: if $f$ is analytic in $|\Im v|<d$, tends to zero there, and $\int|f(v+ib)|dv\le M$ for $|b|<d$, then
\[
\Bigl|h\sum_kf(kh)-\int f\Bigr|\le\frac{2M}{e^{2\pi d/h}-1}.
\]
Take $d=a/2$, inside the poles of $P_a(w-\cdot)$ at $w\pm ia$. On the strip,
\[
|a^2+(t-ib)^2|=|t-i(b+a)|\,|t-i(b-a)|\ge t^2+\tfrac{a^2}4,
\]
so $\int|P_a(w-v-ib)|\,dv\le(a/\pi)(2\pi/a)=2$ and $M\le2\sup|\Psi|$. For
\[
\Psi(v+ib)=\tfrac12\bigl[E(a-b+i(v-c))+E(a+b-i(v-c))\bigr]
\]
write $E(q)=2e^{z^2}-w(iz)$ with $z=q/\sigma\sqrt2$ and $w$ the Faddeeva function, which satisfies $|w(\zeta)|\le1$ for $\Im\zeta\ge0$, that is for $\Re q>0$; and $|e^{z^2}|=e^{((\Re q)^2-(\Im q)^2)/2\sigma^2}\le e^{9a^2/8\sigma^2}$. Hence $S_\Psi$. The decaying bound is Poisson summation: $\varphi_h^\infty(w)-\varphi(w)=\sum_{k\ne0}\widehat{f_w}(2\pi k/h)$ and
\begin{gather*}
\widehat{f_w}(\omega_0)=\frac{e^{-i\omega_0w}}{2\pi}\int\widetilde G(\omega)e^{i\omega(w-c)}d\omega,\\
\widetilde G(\omega)=\sigma\sqrt{2\pi}\,e^{a|\omega|-a|\omega-\omega_0|-\sigma^2\omega^2/2},
\end{gather*}
which is smooth except for kinks at $0$ and $\omega_0$. Two integrations by parts give $(w-c)^{-2}$ times the sum of the jumps of $\widetilde G'$, each at most $2a\sigma\sqrt{2\pi}e^{-a|\omega_0|}$ since $e^{2a|\omega_0|-\sigma^2\omega_0^2/2}\le1$ for $|\omega_0|\ge4a/\sigma^2$, which holds for every sampling frequency $|\omega_0|\ge2\pi/h$ exactly when $h\le\pi\sigma^2/(2a)$; plus $\int|\widetilde G''|\le\sigma\sqrt{2\pi}e^{-a|\omega_0|}J$, using $e^{a|\omega|-a|\omega-\omega_0|}\le e^{-a|\omega_0|+2a|\omega|}$. Summing the geometric series over $k$ gives $E_{\mathrm{dec}}$. For $\psi_1$ the transform carries the factor $-i\sigma^2\omega$, the strip bound is the Fourier integral of $\Psi_1$ on the strip, and the second derivative of $\omega\,\widetilde G$ is bounded in the same way, which gives the constants $S_1$, $C_1$, $J_1$ (derived in the review of the source, not rerun).
\end{proof}

With $h=0.02$ the condition $h\le\pi\sigma^2/(2a)$ holds for all four parameter pairs of the table, $e^{\pi a/h}$ is $e^{78.5}$ at $a=\tfrac12$ and $e^{39.3}$ at $a=\tfrac14$, and the constants are below $10^{-30}$ and $10^{-15}$ respectively and are carried through the computation. The last ingredient is a bound on the mass of $\mathfrak U_{\theta'}$, which is unknown but dominated through the forward equation: for $L>0$,
\[
\begin{gathered}\mathfrak U_{\theta'}([-W,W])\ \le\ \frac{2M_1(W+L)}{1-2P_a((L,\infty))},\\ M_1(T)=\int_0^Tu_1(v)\,dv=\pi^{-1}\Im\log\xi(1+iT),\end{gathered}
\]
since $P_a*\mathfrak U_{\theta'}\ge P_a*(\mathfrak U_{\theta'}|_{[-W,W]})$ and the latter integrates to at least $P_a([-L,L])\,\mathfrak U_{\theta'}([-W,W])$ over $[-W-L,W+L]$. The function $M_1$ is explicit up to $\arg\zeta(1+iT)$ (continuous branches of the logarithms, continued horizontally from $\sigma=+\infty$):
\[
\pi M_1(T)=\arctan T+\tfrac\pi2-\tfrac T2\log\pi+\Im\log\Gamma\bigl(\tfrac{1+iT}2\bigr)+\arg\zeta(1+iT),
\]
and $|\arg\zeta(1+iT)|\le24.303+\log\log T+0.498$ for $T\ge13$, by integrating the explicit bound $|\zeta'/\zeta(\sigma+it)|\le24.303\log t$ on $\sigma\ge1$ from $1$ to $1+1/\log T$, the bound $-\zeta'/\zeta(\sigma)\le1/(\sigma-1)$ from there to $2$, and $\log\zeta(2)=0.4977$ beyond. The explicit bound is Leong's \cite{Leong24}; the constant $24.303$ is the author's, and the corollary of that paper is reported (unchecked) to give $18.062\log t$ for $\sigma\ge1$, $t\ge5$, so $24.303$ is conservative. So $M_1$ is a ball, and the mass of $\mathfrak U_{\theta'}$ on $[-W,W]$ is bounded by an increasing explicit function $\overline U(W)$, with $\overline U(880)=1286.4$ at $a=\tfrac12$ (the author's value, not rerun).

\begin{theorem}[certified windows; conditional, computed by the author]\label{cx:thm:apostcert}
Let $\theta'\in\{\tfrac12,\tfrac34\}$, $a=1-\theta'$, $\sigma\in\{0.3,0.5\}$, suppose $\zeta$ has no zeros in $\Re s>\theta'$ (at $\theta'=\tfrac12$ this is the Riemann hypothesis; at $\theta'=\tfrac34$ it is Theorem~\ref{cx:thm:certified} up to height $1419.9$ and an assumption beyond), and let $\mu$ be the accelerated EM output after $4000$ steps on the grid of spacing $h=0.02$ on $[-440,440]$. For each of the $198$ windows $\varphi$ of width $\sigma$ centred at the grid point $c$ nearest a zero in $(5,395)$, with $\varphi_1=(v-c)\varphi$,
\begin{gather*}
\Bigl|\int\varphi\,d\mu-\int\varphi\,d\mathfrak U_{\theta'}\Bigr|\le B_0(c),\qquad\Bigl|\int\varphi_1\,d\mu-\int\varphi_1\,d\mathfrak U_{\theta'}\Bigr|\le B_1(c),\\
B_0(c)=h\!\!\sum_{|v_i|\le880}\!\!|\psi(v_i)||r_i|+\|\mu\|\,\varepsilon_\mu+\overline U(880)\,\varepsilon_U+T(c),\\ T(c)=\int_{|w|>880}\!\!\!D(|w|-|c|)\,d\mathfrak U_{\theta'}(w),
\end{gather*}
and $B_1$ likewise with $\psi_1$, where $\varepsilon_\mu,\varepsilon_U$ are the uniform bounds on $|\varphi-\varphi_h|$ over the support of $\mu$ and over $[-880,880]$ given by Lemma~\ref{cx:lem:trap} and the tails of Lemma~\ref{cx:lem:psiclosed}, $D$ is the explicit decreasing majorant
\[
D(d)=\frac{E_{\mathrm{dec}}}{d^2}+\frac{K_{\mathrm{far}}\,a}{\pi(1-\lambda)^2d^2}+\Bigl(1+\frac h{\pi a}\Bigr)\bar k(\lambda d),\qquad\lambda=0.9,
\]
with $\bar k$ the tail bound of Lemma~\ref{cx:lem:psiclosed} and $K_{\mathrm{far}}$ a bound on $\sum_{|v_i|>880}|\psi(v_i)|h$, and the last integral is bounded by Abel summation against $\overline U$ over blocks $W_k=880\cdot1.005^k$ up to $10^7$ and by $\overline U(w)\le w\log w$ beyond. Every quantity is a ball. The $792$ pairs of bounds (mass and moment, for $198$ windows, two levels and two widths) were compared with floating point reference values and agree. The certified figures are in the table below.
\end{theorem}

The finite computations are the author's, not rerun; given them, the proof is the identity \eqref{cx:eq:exactid}, Lemma~\ref{cx:lem:trap} and the mass bound. The majorant $D$ is valid after Lemma~\ref{cx:lem:trap}: split the omitted sample points into $|v_i-c|\le\lambda d$, where $P_a(w-v_i)\le a/[\pi(1-\lambda)^2d^2]$, and the complement, where $|\psi(v_i)|\le\bar k(\lambda d)$ and $h\sum_iP_a(w-v_i)\le1+h/(\pi a)$. The Abel summation is an upper sum: if $D$ is decreasing, integration by parts gives $\int_{|w|>V}D(|w|-|c|)\,d\mathfrak U\le\int_V^\infty[-D'(R-|c|)]\,\overline U(R)\,dR$, and on a block $[W_k,W_{k+1}]$ the integral is at most $\overline U(W_{k+1})[D(W_k-|c|)-D(W_{k+1}-|c|)]$.

\begin{center}
\small
\setlength{\tabcolsep}{2.5pt}\footnotesize
\begin{tabular}{llccccccc}
\toprule
$\theta'$ & $\sigma$ & \multicolumn{3}{c}{mass bound $B_0$} & \multicolumn{2}{c}{moment $B_1$} & \multicolumn{2}{c}{centroid bound}\\
& & max & mean & far tail & max & mean & median & max\\
\midrule
$1/2$ & 0.5 & $2.2\times10^{-2}$ & $5.4\times10^{-3}$ & $1.7\times10^{-3}$ & $1.8\times10^{-3}$ & $1.1\times10^{-3}$ & $0.027$ & $0.20$\\
$1/2$ & 0.3 & $1.4\times10^{-2}$ & $5.1\times10^{-3}$ & $1.1\times10^{-3}$ & $2.0\times10^{-3}$ & $1.2\times10^{-3}$ & $0.0072$ & $0.11$\\
$3/4$ & 0.5 & $5.7\times10^{-3}$ & $1.7\times10^{-3}$ & $8.3\times10^{-4}$ & $2.2\times10^{-4}$ & $1.1\times10^{-4}$ & $0.042$ & $0.14$\\
$3/4$ & 0.3 & $3.4\times10^{-3}$ & $1.0\times10^{-3}$ & $4.9\times10^{-4}$ & $3.6\times10^{-5}$ & $1.2\times10^{-5}$ & $0.0068$ & $0.071$\\
\bottomrule
\end{tabular}
\end{center}

\noindent The far tail column is the mean of the last term of $B_0$, the price of knowing $\mathfrak U_{\theta'}$ beyond $880$ only through its mass bound; it is a fifth to a half of the mean bound and would shrink on a longer grid. The main term, the finite sum in \eqref{cx:eq:exactid}, has mean $3.6\times10^{-3}$ at $\theta'=\tfrac12$, $\sigma=0.3$, against $4.5\times10^{-3}$ for the floating point bound with its own tail estimate. At $\theta'=\tfrac12$ every window has certified mass at least $0.9929$, so, since an interval of length one below height $400$ holds at most $n(400)<9.34$ ordinates (the count of Theorem~\ref{cx:thm:fastu1}), each window has a zero within $0.73$ of its centre at $\sigma=0.3$ and within $1.22$ at $\sigma=0.5$: if no zero lay within $r$ of the centre, the Gaussian mass would be at most $2n\sum_{k\ge0}e^{-(r+k)^2/(2\sigma^2)}$, which is $0.967$ and $0.952$ at these radii, below $0.9929$ (the critical radii are $0.7267$ and $1.2115$). The certified centroid lies within $0.05$ of the window centre for $192$ of the $198$ windows at $\sigma=0.3$ (as reported by the author), the six exceptions being the close pairs, where a neighbouring atom at distance $d$ moves the centroid by an amount of the order of $d\,e^{-d^2/2\sigma^2}$. At $\theta'=\tfrac34$ the target is the Thorin density $\sum_\rho P_{3/4-\sigma_\rho}(v-\gamma_\rho)$, which is $\sum P_{1/4}(v-\gamma)$ when the zeros are on the line, the window mass is about $0.7$, the bounds are five times tighter because the multiplier $e^{|\omega|/4}$ is milder. The serialisation of the EM vector and of the data balls must be lossless for the bound to be a certificate; this was not repeated here.

\begin{corollary}[from centroids to zeros]\label{cx:cor:zeroloc}
Assume the Riemann hypothesis, so that $\mathfrak U_{1/2}=\sum_\rho\delta_{\gamma_\rho}$. Let $\varphi(v)=e^{-(v-c)^2/(2\sigma^2)}$, let $\int\varphi\,d\mathfrak U_{1/2}\in[\widehat M-B_0,\widehat M+B_0]$ with $\widehat M>B_0$ and $\int(v-c)\varphi\,d\mathfrak U_{1/2}\in[\widehat m-B_1,\widehat m+B_1]$ be certified, let $\gamma_*$ be an ordinate with $|\gamma_*-c|\le R$, and let $\Theta\ge\sum_{\gamma\ne\gamma_*}(|\gamma-c|+R)\varphi(\gamma)$, the sum over the other ordinates with multiplicity. Then
\[
|\gamma_*-c|\ \le\ \frac{|\widehat m|+B_1+\Theta}{\widehat M-B_0}.
\]
\end{corollary}

\begin{proof}
With $M=\int\varphi\,d\mathfrak U_{1/2}$, $\int(v-c)\varphi\,d\mathfrak U_{1/2}=(\gamma_*-c)M+\sum_{\gamma\ne\gamma_*}(\gamma-\gamma_*)\varphi(\gamma)$ and $|\gamma-\gamma_*|\le|\gamma-c|+R$. Hence $|\gamma_*-c|\,M\le|\widehat m|+B_1+\Theta$, and $M\ge\widehat M-B_0>0$.
\end{proof}

The windows are centred at grid points, within $0.01$ of the zeros, so that $\psi(v_i-c)$ is a shift of one table; the centroid certificate reports the offset from $c$ and is unaffected. Below height $1419.9$ the ordinates are already certified (Theorem~\ref{cx:thm:certified}), so Corollary~\ref{cx:cor:zeroloc} locates nothing new there. The hypothesis on the zeros is used only to know that $\mathfrak U_{\theta'}$ exists and satisfies $P_a*\mathfrak U_{\theta'}=u_1$. The bound is a statement about $\mathfrak U_{\theta'}$: at $\theta'=\tfrac34$ about the Thorin density itself, at $\theta'=\tfrac12$, under the Riemann hypothesis, about Gaussian weighted sums over the zeros, and about single zeros only through Corollary~\ref{cx:cor:zeroloc}.

\subsection{An iteration adapted to the certificate}\label{cx:sec:weighted}

The certificate's main term $h\sum_i|\psi_c(v_i)||r_i|$ is a weighted $\ell^1$ norm of the residual; does an iteration aimed at it beat EM? We compared, at $\theta'=\tfrac12$ on the grid of spacing $0.02$, with $\sigma=0.3$ and the weights $w_i=\sum_c|\psi_c(v_i)|$ summed over the $198$ windows: accelerated EM; FISTA with nonnegativity on the weighted least squares $\tfrac12\sum_iw_i(Ax-y)_i^2$, which bounds the sum over the windows of the weighted $\ell^1$ norms by Cauchy and Schwarz, $h\sum_iw_i|r_i|\le h(\sum_iw_i)^{1/2}(\sum_iw_ir_i^2)^{1/2}$ (an aggregate bound, not a bound for each window); FISTA on the plain least squares; and accelerated EM followed by a weighted FISTA polish. The figure of merit is the mean main term over the windows. The experiment is the author's, not repeated.

\begin{center}
\small
\begin{tabular}{lcc}
\toprule
iteration (steps) & main term & KL\\
\midrule
accelerated EM (4000) & $5.32\times10^{-5}$ & $1.2\times10^{-4}$\\
weighted FISTA (4000) & $7.38\times10^{-5}$ & $0.95$\\
plain FISTA (4000) & $5.26\times10^{-5}$ & $1.4\times10^{-4}$\\
EM (4000), weighted FISTA polish (1000) & $4.12\times10^{-5}$ & $4.3\times10^{-4}$\\
accelerated EM (5000) & $4.73\times10^{-5}$ & $9.8\times10^{-5}$\\
\bottomrule
\end{tabular}
\end{center}

\noindent The weighted objective from a cold start is worse on its own figure of merit, and much worse in divergence: the weights concentrate on the peaks of $u_1$, where the deconvolution is most ill conditioned, and the iteration spends its steps there. Plain least squares and EM are indistinguishable. The polish gains thirteen per cent over the same number of EM steps. Our reading is that the main term is dominated by the residual within a unit of each zero, where every objective fits equally badly, the limit being the conditioning rather than the loss. In this experiment the certificate does not reward an iteration of its own; a tighter certificate needs more accelerated steps or a finer grid near the zeros. One experiment does not show that no other objective helps.

\subsection{A fast algorithm for $u_1$ at fixed accuracy}\label{cx:sec:fastu1}

The third problem asked for an algorithm for $u_1(v)$ polynomial in $\log v$. With the accuracy entering through $\log(1/\varepsilon)$ it is open (Problem~\ref{cx:prob:polylog}). With the accuracy fixed the explicit formula answers it: given that the zeros near the height $v$ lie on the critical line (verified below $3\times10^{12}$ \cite{PT21}; at every height conditionally on the announced half plane $\Re s>\tfrac78$, with a larger exponent), the primes below $X$ determine $u_1$ at that height to accuracy of order $X^{-1/2}\log v$, so the cost of one value is $O(X)$ whatever the height, and a nonuniform fast Fourier transform gives a whole grid of heights for the price of one.

\begin{theorem}[smoothed prime sum]\label{cx:thm:fastu1}
Let $v\ge2$, $X\ge2$, $s=1+iv$, $L=\zeta'/\zeta$, and
\begin{multline*}
u_1^X(v)=\frac1\pi\Bigl[\frac1{1+v^2}-\frac{\log\pi}2+\frac12\Re\psi\Bigl(\frac{1+iv}2\Bigr)\\-\sum_{n\ge2}\frac{\Lambda(n)}n\,e^{-n/X}\cos(v\log n)\Bigr].
\end{multline*}
Then, with $R_X(s)=(2\pi i)^{-1}\int_{\Re w=-3/2}(-L(s+w))\Gamma(w)X^w\,dw$,
\[
\begin{split}\pi\bigl(u_1(v)-u_1^X(v)\bigr)&=-\Re\sum_\rho\Gamma(\rho-s)X^{\rho-s}+\Re\,\Gamma(1-s)X^{1-s}\\&\quad+\Re\bigl[X^{-1}L(s-1)\bigr]+\Re R_X(s),\end{split}
\]
with $|R_X(s)|\le2.5X^{-3/2}\log(v+2)$, and if every zero with $|\gamma-v|\le Y$, $Y\ge1$, lies on the critical line,
\[
\begin{split}|u_1(v)-u_1^X(v)|&\le\frac1\pi\Bigl[c_\Gamma\,n(v)\,X^{-1/2}+2.6\,c_\Gamma'\,n(v+Y+1)\,e^{-\pi Y/2}\\&\qquad+\sqrt{\frac{\pi}{v\sinh\pi v}}+\frac{27\log v}X+\frac{2.5\log(v+2)}{X^{3/2}}\Bigr],\end{split}
\]
where $n(t)=0.39\log t+7$ bounds the number of ordinates in any interval of length one below height $t$, by \cite[(1.7), (2.7) and Corollary~1.4]{HSW22}; $c_\Gamma=\sup_{v\ge2}\sum_{j\in\mathbb Z}\sup_{j\le y<j+1}|\Gamma(-\tfrac12+iy)|\,\frac{n(v+|j|+1)}{n(v)}\le8.34$ (without the factor $n(v+|j|+1)/n(v)$ the sum is $8.14376$); and $c_\Gamma'=\sup\{|\Gamma(x+iy)|e^{\pi|y|/2}:-1<x<0,\ |y|\ge1\}\le2.51$.
\end{theorem}

\begin{proof}
$\sum_n\Lambda(n)n^{-s}e^{-n/X}=(2\pi i)^{-1}\int_{(2)}(-L)(s+w)\Gamma(w)X^w\,dw$ by Mellin inversion of $e^{-y}$. Shift the contour to $\Re w=-\tfrac32$. This crosses the pole of $\Gamma$ at $w=0$, with residue $-L(s)$, the zeros $w=\rho-s$ (real part in $(-1,0)$), with residues $-\Gamma(\rho-s)X^{\rho-s}$ times the multiplicity, the pole $w=1-s$ of $L$, with residue $\Gamma(1-s)X^{1-s}$, and the pole $w=-1$ of $\Gamma$, with residue $X^{-1}L(s-1)$; it crosses none of the trivial zeros ($\Re w=-2k-1\le-3$) or of the later poles of $\Gamma$. The horizontal segments vanish along heights that avoid the zeros, by the exponential decay of $\Gamma$. Taking real parts and using $\pi u_1(v)=\Re[1/s-\tfrac12\log\pi+\tfrac12\psi(s/2)]+\Re L(s)$ gives the identity. On $\Re w=-\tfrac32$, $\Re(s+w)=-\tfrac12$; by the functional equation $L(z)=\chi'/\chi(z)-L(1-z)$ with $\Re(1-z)=\tfrac32$, where the Dirichlet series converges absolutely, so $|L(1-z)|\le-\zeta'/\zeta(\tfrac32)<1.5053$. Here $\chi'/\chi(z)=\log2\pi+\tfrac\pi2\cot\tfrac{\pi z}2-\psi(1-z)$ with $|\cot\tfrac{\pi z}2|=1$ on $\Re z=-\tfrac12$, and Binet's formula $\psi(a)=\log a-\tfrac1{2a}-2\int_0^\infty\frac{t\,dt}{(t^2+a^2)(e^{2\pi t}-1)}$, with $|t^2+a^2|\ge\tfrac94$ for $\Re a=\tfrac32$, gives $|\psi(\tfrac32+i\tau)|\le\log(|\tau|+2)+\tfrac\pi2+\tfrac13+\tfrac1{27}$. For $w=-\tfrac32+iy$, as $\log(|v+y|+2)\le\log(v+2)+\log(1+|y|)$, this gives $|L(s+w)|\le\log(v+2)+\log(1+|y|)+6.856$, and with $\int_{\mathbb R}|\Gamma(-\tfrac32+iy)|\,dy<2.527$ and $\int_{\mathbb R}|\Gamma(-\tfrac32+iy)|\log(1+|y|)\,dy<0.798$ (computed), $|R_X(s)|\le X^{-3/2}\bigl(0.41\log(v+2)+2.89\bigr)\le2.5X^{-3/2}\log(v+2)$, since $\log(v+2)\ge\log4$.

For the bound: the pole term is $\Re\Gamma(-iv)X^{-iv}$, of modulus at most $|\Gamma(-iv)|=(\pi/(v\sinh\pi v))^{1/2}$ since $|X^{-iv}|=1$. For $v\ge5$, $|L(iv)|\le27\log v$ by the functional equation, $|\chi'/\chi(iv)|\le\log v$ and the explicit bound of \cite{Leong24} for $|L(1-iv)|$ on $\Re s=1$; for $2\le v\le5$ it is a direct evaluation (the maximum of $|L(iv)|/\log v$ on $[2,60]$ is $1.014$). For the zeros on the line, $|\Gamma(-\tfrac12+iy)|^2=\pi/(\cosh(\pi y)(\tfrac14+y^2))$ and $|X^{\rho-s}|=X^{-1/2}$; the zeros are counted in the unit intervals $\gamma-v\in[j,j+1)$, each holding at most $n(v+|j|+1)$ ordinates, and the resulting sum is $c_\Gamma n(v)X^{-1/2}$. A zero at distance more than $Y$ contributes at most $c_\Gamma'e^{-\pi|\gamma-v|/2}$ ($|X^{\rho-s}|\le1$ and $|\Gamma(\sigma_\rho-1+i(\gamma-v))|\le c_\Gamma'e^{-\pi|\gamma-v|/2}$), and $\sum_{k\ge0}2e^{-\pi k/2}n(v+Y+k+1)\le2.6\,n(v+Y+1)$. The constants $8.3388$, $8.14376$ and $2.509$, and the count $n(t)$ against \cite{HSW22}, were recomputed.
\end{proof}

The pole term cannot be dropped. At $v=2$ every zero has $|\gamma-2|\ge12$ and lies on the critical line below $3\times10^{12}$; taking $Y=10^{12}$, the bound without the pole term tends to $0$ as $X\to\infty$. But $\Re\Gamma(-2i)X^{-2i}=|\Gamma(-2i)|\cos(\arg\Gamma(-2i)-2\log X)$ equals $|\Gamma(-2i)|=0.0766$ at $X_k=\exp((\arg\Gamma(-2i)+2\pi k)/2)$, and the other terms tend to $0$, so $\limsup_{X\to\infty}|u_1(2)-u_1^X(2)|\ge0.0766/\pi=0.0244$. At large heights the pole term is exponentially small and does not affect the experiments below.

The hypothesis on the zeros within $Y$ of $v$ is verified for $v+Y\le3\times10^{12}$ with $Y$ of a few dozen \cite{PT21}, and $e^{-\pi Y/2}$ is then negligible. Beyond, under the zero free region alone, $|X^{\rho-s}|\le X^{-c/\log v}$ and the first term decays only when $\log X\gg\log v$. Conditionally on the announced half plane $\Re s>\tfrac78$ \cite{OAI26a,OAI26b,Liu26a}, every zero has $\tfrac18\le\beta\le\tfrac78$, so $|X^{\rho-s}|\le X^{-1/8}$ and $\Re(\rho-s)$ stays in $[-\tfrac78,-\tfrac18]$, away from the poles of $\Gamma$; the same count gives $|u_1(v)-u_1^X(v)-\pi^{-1}\Re\Gamma(1-s)X^{1-s}|\le C\log(v+2)(X^{-1/8}+X^{-1})$ at every height, so that a fixed accuracy $\varepsilon$ costs $X$ of order $(\log v/\varepsilon)^8$. The sum can be cut at $n=30X$ with a dropped tail of about $e^{-30}\log X$ (at about $X\log(1/\delta)$ for a tolerance $\delta$), so one value costs $O(X)$ operations after a sieve of $O(X\log\log X)$, independent of $v$; the bit cost adds the phases $v\log n$, which need $\log n$ to about $\log(v/\varepsilon)$ bits. For accuracy $\varepsilon$ take $X=(c_\Gamma n(v)/\varepsilon)^2$: the cost is $O(\varepsilon^{-2}\log^2v)$ operations, polynomial in $\log v$ and in $1/\varepsilon$.

On a grid $v_k=kh$ the sum is
\[
\Re\sum_na_ne^{ik(h\log n)},\qquad a_n=\frac{\Lambda(n)}n\,e^{-n/X},
\]
a type one nonuniform discrete Fourier transform with the frequencies $h\log n$, and a nonuniform FFT evaluates it at $K$ heights with $M$ retained prime powers in $O(K\log K+M\log(1/\delta))$ operations for a transform accuracy $\delta$ \cite{FINUFFT19}, to which the coefficients, the sieve and the smooth gamma factor part are to be added. The experiment is the author's, not rerun. With $X=10^6$ and $1.86$ million prime powers below $3\times10^7$, the error $|u_1-u_1^X|$ at the heights $10^3$, $10^4$, $10^5$, $10^6$ and $10^7$ is reported as $9\times10^{-4}$, $4\times10^{-4}$, $8\times10^{-4}$, $1.1\times10^{-3}$ and $1.4\times10^{-3}$, against Arb \cite{Arb}; averaged over fifty one heights near $10^6$ the error is $0.05\log v\,X^{-1/2}$ for $X=10^4$, $10^5$ and $10^6$, which is the law of the theorem with a constant about twenty times smaller than the leading coefficient of the proved bound ($c_\Gamma\cdot0.39/\pi=1.0$ per $\log v$), the zero sum being oscillatory. The nonuniform FFT (\texttt{finufft}) is reported to give $u_1^X$ at $1{,}048{,}576$ heights at spacing $0.02$, up to height $20{,}972$, in $0.4$ seconds for the prime sums alone (the smooth part is not in the timing and the thread count is not fixed), agreeing with the direct sums to $10^{-8}$; Arb needs about two milliseconds per height, two thousand seconds for the same grid, but gives twenty digits.

So for the data layer: twenty digits at one height cost of order $v^{1/2}$ by the formula of Riemann and Siegel \cite{Edwards,Siegel32} and of order $v^{1/3}$ by Hiary \cite{Hiary11}; three digits at a million heights cost a second by the primes; the deconvolution needs four or five digits on a grid, which by the observed law above the smoothed sum delivers with $X=10^8$ to $10^{10}$ (the proved bound asks for more), at the edge of a laptop, and Arb delivers in minutes up to the heights of this chapter.

\subsection{Is $\mathfrak U_{1/2}$ computed efficiently?}\label{cx:sec:halfeff}

At the bottom level the deconvolution $P_{1/2}*\mathfrak U_{1/2}=u_1$ is, in Fourier variables, the explicit formula read backwards: the multiplier $e^{|\omega|/2}$ turns the atoms $-\Lambda(n)/n$ of $\widehat{u_1}$ into the weights $-\Lambda(n)/\sqrt n$ of Weil's formula (Theorem~\ref{es:thm:explicit}). The heuristic answer on efficiency is Remark~\ref{es:rem:resolution}: to resolve $\mathfrak U_{1/2}$ at scale $\delta$ needs the primes to $e^{\pi/\delta}$, or data accurate to $e^{-\pi/2\delta}$, so the problem is exponentially ill posed in the resolution and no linear band limited inversion of $u_1$ escapes this; the EM computation reached $\delta\approx10^{-3}$ only because its target is not a general measure (Section~\ref{es:sec:level}). This is worst case stability, not a lower bound for every method: for a single unit atom with blur $y(v)=P_a(v-\gamma)$, $\gamma=a\,[1/y(-t)-1/y(t)]/(4\pi t)$ for any $t\ne0$, so extra structure changes the inverse problem. The backward flow with $4000$ accelerated steps on a grid of spacing $0.02$ certifies Gaussian weighted centroids with an uncertainty of about $10^{-3}$ (the moment bound $B_1$ of Theorem~\ref{cx:thm:apostcert}, divided by the mass), not the atoms (Corollary~\ref{cx:cor:zeroloc}), and further digits cost further steps, at the observed rates of Section~\ref{cx:sec:rateworked}, or a finer grid with data accurate to about $e^{-\pi/2h}$. Each FFT step costs $O(T\log T)$ on a grid to height $T$, and the number of steps for a fixed resolution was observed to be grid independent; no bound on it uniform in $T$, the least gap or the accuracy is proved, so the total $T^{1+\varepsilon}$ of Odlyzko and Sch\"onhage \cite{OdlyzkoSchonhage88} is not claimed for the deconvolution. In the sense that matters it is not efficient: direct zero finding gives $\gamma$ far more accurately in less time, because the functional equation carries the spectrum beyond any $\Omega$ in closed form, where the deconvolution has only the data below $\Omega$. The efficient algorithm for $\mathfrak U_{1/2}$ is Turing's \cite{Turing53}; the deconvolution contributes the certificate, a bound on the measure from the residual alone, and the levels above $\tfrac12$, where the Thorin measure is a density that no zero finder produces.

\subsection{A von Neumann representation of $\mathfrak U_{1/2}$}\label{cx:sec:vonneumann}

Two things could be meant. In the sense of the Neumann series of Section~\ref{cx:sec:vnmc} (the Richardson iteration with unit step, whose symbol $1-e^{-|\omega|/2}$ lies in $[0,1)$), the sum to $K$ terms applied to $u_1$ has symbol $[1-(1-e^{-|\omega|/2})^{K+1}]e^{|\omega|/2}$, so relative to the target the effective filter is $1-(1-e^{-|\omega|/2})^{K+1}$, which is near $1$ for $|\omega|\le2\log K$ and near $0$ beyond. The truncated Neumann series is therefore a soft band limited inverse with resolution of order $\pi/(2\log K)$, and the resolution $0.1$ costs $K=e^{16}$ terms. That representation is exact, convergent and useless, for the reason of the previous subsection.

In the sense of the spectral theorem, the question is whether $\mathfrak U_{1/2}$ is the spectral measure of a natural self adjoint operator, the conjecture of Hilbert and P\'olya. Under the Riemann hypothesis a diagonal operator with eigenvalues $\gamma_\rho$ exists trivially, so existence alone is not the conjecture. What is unconditional is the shape of the answer at the levels above. If $H$ is self adjoint with spectral measure $\mu_H=\sum\delta_{\gamma}$, then for $a>0$ the resolvent is not trace class, since $N(T)\asymp T\log T$, but its imaginary part is, and
\[
\frac1\pi\operatorname{Tr}\,\Im(H-v-ia)^{-1}=\sum_\gamma\frac{a/\pi}{(\gamma-v)^2+a^2}=(P_a*\mu_H)(v),
\]
so the Thorin density at level $\theta=\tfrac12+a$ is this resolvent trace at height $a$, and the backward flow $\theta\downarrow\tfrac12$ is Stieltjes inversion. Unconditionally the chapter has computed the function that would be that trace:
\[
u_\theta(v)=\frac1\pi\Re\frac{\xi'}\xi(\theta+iv)=\frac1\pi\sum_\rho\Re\frac1{\theta+iv-\rho}=\frac1\pi\sum_\rho\frac{a_\rho}{a_\rho^2+(v-\gamma_\rho)^2},
\]
with $a_\rho=\theta-\sigma_\rho$, which is that trace exactly when all $a_\rho$ equal $a$. So the densities $u_\theta$, $\theta>\tfrac12$, are the traces $\pi^{-1}\operatorname{Tr}\Im(H-v-i(\theta-\tfrac12))^{-1}$ of the one operator $H$ with spectrum $\{\gamma_\rho\}$ if and only if the Riemann hypothesis holds. (If a zero has $\sigma_0>\tfrac12$, then at $v=\gamma_0$ the density $u_\theta$ is unbounded as $\theta\downarrow\sigma_0$, while the trace of $\Im$ of the resolvent is finite for every $\theta>\tfrac12$.) The Thorin measures at levels $\theta>\tfrac12$ provide the resolvent trace without the operator: $\xi'/\xi$ plays the trace of the resolvent at $z=v+i(\theta-\tfrac12)$, and Theorem~\ref{cx:thm:apostcert} bounds the boundary measure from the resolvent at height $\tfrac12$, the Stieltjes inversion made quantitative. This does not produce $H$; it says what $H$ would have to do.

\section{Open problems}\label{cx:sec:open}

\begin{problem}\label{cx:prob:twophase}
Prove the two phase picture of Section~\ref{cx:sec:rateworked} for the continuous Cauchy deconvolution at an atomic target: a sublinear phase with $k(R_k-R^*)$ near a constant of order $50$, lasting of the order of $10^4$ steps on this data, followed by faster convergence on the support of the atoms with a grid independent rate; and compute that rate from the curvature of the divergence at the atoms. For the accelerated scheme with restart, explain the observed three point exponents, which oscillate between $0.6$ and $2.8$ with a running value near $1.3$.
\end{problem}

\begin{problem}\label{cx:prob:closedform}
Proposition~\ref{cx:prop:slope} gives the limiting slope under the Riemann hypothesis to six digits, $0.183899$ at $v^*=17.4979$. Is there a closed form? By the Hadamard product, under the Riemann hypothesis, $\log|\xi(\tfrac12+iv)|=\log\xi(\tfrac12)+\sum_{\gamma>0}\log|1-v^2/\gamma^2|$, so with $\mathcal C(v)=\partial_v^2\log|\xi(\tfrac12+iv)|$ the bracket is $X(v)=\mathcal C(0)-\mathcal C(v)$: the slope of the excess is the drop of the second logarithmic derivative of $\xi$ along the critical line from the origin to $v$, and $v^*$ is the zero of $\mathcal C'$ between the first two zeros. Whether $\mathcal C(v^*)$ has an expression in known constants is not clear.
\end{problem}

\begin{problem}\label{cx:prob:polylog}
Theorem~\ref{cx:thm:fastu1} evaluates $u_1$ to accuracy $\varepsilon$ in time polynomial in $\log v$ and in $1/\varepsilon$. Is there an algorithm polynomial in $\log v$ and $\log(1/\varepsilon)$? Such an algorithm for $\zeta$ and $\zeta'$ on $\Re s=1$ would give one; the converse is not known, and the question is open. The smoothed sum cannot be the route, since the zeros themselves contribute $X^{-1/2}$.
\end{problem}

\begin{problem}\label{cx:prob:bestresidual}
The certificate's main term is set by the residual within a unit of each zero, which no objective reduces faster than the conditioning allows (as observed in Section~\ref{cx:sec:weighted}). Quantify this: for a single atom at $\gamma$ blurred by $P_a$ and observed exactly, what is the best achievable $h\sum_i|\psi_\gamma(v_i)||r_i|$ after $k$ steps of any first order method on the cone, as a function of $k$, $h$ and $a$? The answer would say how far Theorem~\ref{cx:thm:apostcert} can be pushed at a given cost.
\end{problem}
